\newcolumntype{Y}[1]{>{\raggedright\arraybackslash\hspace{0pt}}p{#1}}

\begin{table}[H]
    \centering
    \caption[研究假設、比較對象與評估指標]{研究假設、比較對象與評估指標}\label{table:hypotheses}
    \renewcommand{\arraystretch}{1.3}
    \setlength{\tabcolsep}{5pt}
    \normalsize
    \begin{tabular}{|Y{\dimexpr 0.1200\textwidth-2\tabcolsep\relax}|Y{\dimexpr 0.3800\textwidth-2\tabcolsep\relax}|Y{\dimexpr 0.5000\textwidth-2\tabcolsep\relax}|}
        \hline
        \textbf{假設} & \textbf{比較對象} & \textbf{看的指標} \\ \hline
        H1 & 關鍵字找餐廳、向量、混合三種召回 & 排名指標（於先導階段確定） \\ \hline
        H2a & 完整證據驗證 vs. AI 直接排名 & 違反限制率、未知誤承諾率；同時報告覆蓋率與拒答率 \\ \hline
        H2b（探索） & 圖譜 vs. 關聯式資料庫版本 & 查詢正確性與證據完整度（指標定義於先導階段確定） \\ \hline
        H3 & 以預估抵達時間檢查營業時間 vs. 以出發時間檢查 & 營業違反率（到店時已打烊的比例），依用餐時段分層報告 \\ \hline
        H4 & 加入即時到站查核 vs. 只用旅運規劃回傳的時間 & 無法上車或轉乘接不上的比例；「不確定」標示的正確率（以現場驗證為標準） \\ \hline
        H5（探索） & 餐廳與搭乘方案合併排序 vs. 先選餐廳再配公車 & 最終方案的總旅程時間、等車時間與違反率；排名因搭乘方案而改變的比例 \\ \hline
        H6 & 行程中監控並重新規劃 vs. 不監控的靜態方案 & 偵測率、誤報率、提醒領先時間與替代方案可搭乘比例（於重播資料上計算，現場結果僅描述） \\ \hline
    \end{tabular}
\end{table}
